\documentclass[aspectratio=169]{beamer}

\usetheme{Madrid}
\usecolortheme{seahorse}
\usefonttheme{professionalfonts}

\usepackage{listings}
\usepackage{xcolor}
\usepackage{tikz}
\usepackage{booktabs}
\usepackage{hyperref}
\usepackage{textcomp}
\usepackage{microtype}

\definecolor{codebg}{HTML}{1E1E2E}
\definecolor{codegreen}{HTML}{A6E3A1}
\definecolor{codeblue}{HTML}{89DCEB}
\definecolor{codeyellow}{HTML}{F9E2AF}
\definecolor{codegray}{HTML}{6C7086}
\definecolor{codepurple}{HTML}{CBA6F7}
\definecolor{codeorange}{HTML}{FAB387}
\definecolor{codewhite}{HTML}{CDD6F4}
\definecolor{alertred}{HTML}{F38BA8}

\lstdefinestyle{clang}{
  backgroundcolor=\color{codebg},
  basicstyle=\ttfamily\footnotesize\color{codewhite},
  keywordstyle=\color{codepurple}\bfseries,
  commentstyle=\color{codegray}\itshape,
  stringstyle=\color{codegreen},
  numberstyle=\tiny\color{codegray},
  numbers=left, numbersep=6pt,
  breaklines=true, frame=single,
  rulecolor=\color{codeblue},
  language=C,
  tabsize=4
}

\title[MyOS Chapter 3]{\textbf{MyOS Chapter 3: Full Keyboard I/O + Scrollback}}
\subtitle{From basic IRQ keys to interactive terminal navigation}
\author{MyOS Project}
\date{March 2026}
\institute{Interrupt-driven input subsystem}

\newcommand{\concept}[2]{%
  \begin{block}{#1}\smallskip #2\smallskip\end{block}\vspace{5pt}}

\newcommand{\warn}[1]{%
  \begin{alertblock}{Important}\smallskip #1\smallskip\end{alertblock}\vspace{5pt}}

\begin{document}

\begin{frame}
  \titlepage
\end{frame}

\begin{frame}{Why scrolling failed before}
  \begin{itemize}
    \item Previous terminal driver stored only visible 80x25 VGA cells
    \item Once line 26 arrived, line 1 was physically overwritten
    \item Arrow keys had no viewport concept to move through old content
  \end{itemize}
  \vspace{6pt}
  \concept{Root cause}{
    No persistent text history buffer existed in RAM. VGA memory was the only source of truth.
  }
\end{frame}

\begin{frame}[fragile]{Scrollback Architecture in \texttt{screen.c}}
  \begin{lstlisting}[style=clang]
#define SCROLLBACK_LINES 512

static uint16_t text_buffer[SCROLLBACK_LINES][VGA_WIDTH];
static int cursor_row = 0;
static int cursor_col = 0;
static int total_lines = 1;
static int view_top_line = 0;
  \end{lstlisting}
  \textbf{Logic behind these lines:}
  \begin{itemize}
    \item \texttt{text\_buffer} keeps 512 rows of terminal history in RAM
    \item \texttt{cursor\_row/col} track absolute write position in history
    \item \texttt{view\_top\_line} tracks which history row is shown at VGA row 0
    \item VGA now becomes a rendered viewport, not the primary storage
  \end{itemize}
\end{frame}

\begin{frame}[fragile]{Render Pipeline: History \textrightarrow VGA}
  \begin{lstlisting}[style=clang]
static void render_view() {
    for (int row = 0; row < VGA_HEIGHT; row++) {
        int source_row = view_top_line + row;
        for (int col = 0; col < VGA_WIDTH; col++) {
            if (source_row < total_lines)
                video_memory[row * VGA_WIDTH + col] =
                    text_buffer[source_row][col];
            else
                video_memory[row * VGA_WIDTH + col] = make_cell(' ');
        }
    }
}
  \end{lstlisting}
  \concept{Rendering model}{
    Every key/output updates history first, then \texttt{render\_view()} redraws the 25-line window.
  }
\end{frame}

\begin{frame}[fragile]{Viewport Controls}
  \begin{lstlisting}[style=clang]
void scroll_view_up(int lines) {
    if (view_top_line >= lines) view_top_line -= lines;
    else view_top_line = 0;
    render_view();
    update_cursor();
}

void scroll_view_down(int lines) {
    int max_top = total_lines > VGA_HEIGHT ? total_lines - VGA_HEIGHT : 0;
    view_top_line += lines;
    if (view_top_line > max_top) view_top_line = max_top;
    render_view();
    update_cursor();
}
  \end{lstlisting}
  \textbf{Result:} Up/Down keys can navigate historical output safely.
\end{frame}

\begin{frame}[fragile]{Keyboard IRQ: Full Scancode Handling}
  \begin{lstlisting}[style=clang]
if (scancode == 0xE0) { extended_prefix = 1; ... }

if (extended_prefix) {
    extended_prefix = 0;
    if (scancode == 0x48) scroll_view_up(1);   // Arrow Up
    else if (scancode == 0x50) scroll_view_down(1); // Arrow Down
    ...
}

if (scancode == 0x2A || scancode == 0x36) shift_pressed = 1;
if (scancode == 0xAA || scancode == 0xB6) shift_pressed = 0;
if (scancode == 0x3A) caps_lock = !caps_lock;
  \end{lstlisting}
  \begin{itemize}
    \item Handles Shift make/break events, Caps lock toggle
    \item Handles extended scancode prefix for navigation keys
    \item Filters key-release events before ASCII translation
  \end{itemize}
\end{frame}

\begin{frame}[fragile]{Keymap Translation (Normal + Shifted)}
  \begin{lstlisting}[style=clang]
static const char keymap[128] = { ... '1','2','3',...,'a','b',...,'/' ... };
static const char keymap_shift[128] = { ... '!','@','#',...,'A','B',...,'?' ... };

if (is_alpha(normal)) {
    c = (shift_pressed ^ caps_lock) ? shifted : normal;
} else {
    c = shift_pressed ? shifted : normal;
}
  \end{lstlisting}
  \concept{Why XOR for letters?}{
    Shift and Caps both invert case. Pressing both returns to lowercase.
  }
\end{frame}

\begin{frame}[fragile]{Expanded low-level I/O API}
  \begin{lstlisting}[style=clang]
uint8_t  inb(uint16_t port);   void outb(uint16_t port, uint8_t value);
uint16_t inw(uint16_t port);   void outw(uint16_t port, uint16_t value);
uint32_t inl(uint16_t port);   void outl(uint16_t port, uint32_t value);
void io_wait();
  \end{lstlisting}
  \textbf{Why this matters:}
  \begin{itemize}
    \item Keeps kernel hardware access centralized and reusable
    \item Prepares for timer/PIT, ATA, PCI, and future drivers
  \end{itemize}
\end{frame}

\begin{frame}[fragile]{\texttt{print\_int()} reliability fix}
  \begin{lstlisting}[style=clang]
if (n < 0) {
    put_char('-');
    value = (unsigned int)(-(n + 1)) + 1;
} else {
    value = (unsigned int)n;
}
  \end{lstlisting}
  \concept{Why this change}{
    Avoids overflow for INT\_MIN while still printing signed integers correctly.
  }
\end{frame}

\begin{frame}{Testing checklist in QEMU}
  \begin{enumerate}
    \item Type letters/digits/symbols (verify full keymap)
    \item Press Shift + digits (expect \texttt{!@\#...})
    \item Toggle Caps Lock and verify letter case logic
    \item Use Backspace, Tab, Enter
    \item Generate many lines, then Arrow Up/Down for scrollback
  \end{enumerate}
  \warn{If viewport is at history top and new output arrives, terminal follows bottom only when already at bottom before writing.}
\end{frame}

\begin{frame}{What this chapter achieved}
  \begin{itemize}
    \item Full interrupt-driven keyboard input translation
    \item Terminal history buffer with viewport navigation
    \item Cleaner low-level I/O surface for future device drivers
    \item More robust integer printing path for diagnostics
  \end{itemize}
  \vspace{8pt}
  \begin{center}
    \Large \textbf{MyOS is now an interactive, navigable console}
  \end{center}
\end{frame}

\end{document}
